\subsubsection{The local five-dimensional quotient with two three-dimensional modules}

\begin{lemma}\label{B:top-heavy-Tor}
Let $A$ be a finite-dimensional local algebra over a field $k$, with
$A/\operatorname{rad}A=k$.  Let $U$ be a finite-dimensional left $A$-module,
and put
\[
 a=\dim_k(U/JU),\qquad b=\dim_k JU,
 \qquad J=\operatorname{rad}A.
\]
If $a>b$ and $W$ is a finite-dimensional right $A$-module satisfying
$\operatorname{Tor}^A_i(W,U)=0$ for all sufficiently large $i$, then $W$
is projective.  The same assertion holds with left and right interchanged.
\end{lemma}
\begin{proof}
Take a minimal free right resolution of $W$, with term $A^{\beta_i}$ in
degree $i$.  Every differential has entries in $J$.  After tensoring with
$U$, its image in degree $i-1$ lies in $(JU)^{\beta_{i-1}}$.  Consequently
\[
 \dim_k\operatorname{Tor}^A_i(W,U)
 \geq (a+b)\beta_i-b\beta_{i-1}-b\beta_i
 =a\beta_i-b\beta_{i-1}.
\]
Vanishing in all sufficiently large degrees forces
$\beta_i\leq(b/a)\beta_{i-1}$.  Since $b/a<1$, the nonnegative integer
sequence $\beta_i$ must eventually be zero.  A finite-dimensional module
of finite projective dimension over a finite-dimensional local algebra
is free: in a nontrivial finite minimal free resolution, the last
nonzero differential, whose entries lie in $J$, annihilates a nonzero
vector with entries in the appropriate socle of $A$, contradicting its
injectivity.  Thus $W$ is projective.
\end{proof}

\begin{proposition}\label{B:local-quotient-CP2}
Let $M$ be a finite monoid whose minimal ideal $K$ is a rectangular
band.  If $A=\mathbb Q M/\mathbb Q K$ is local and
$\dim_{\mathbb Q}H^2(BM;\mathbb Q)=1$, then every rational
degree-two class has square zero.  In particular $BM$ is not
homotopy equivalent to $\mathbb{CP}^2$.  This assertion has no bound
on the order of $M$ and does not require finite homotopy type.
\end{proposition}
\begin{proof}
Let $W,V$ be the column and row augmentation modules in
Lemma~\ref{lem:B-band-Tor}, and set $J=\operatorname{rad}A$ and
$D=A/J$.  There is a canonical surjection
\[
 W\otimes_A V\longrightarrow
 (W/WJ)\otimes_D(V/JV).
\]
The source has rational dimension one.  Both tops are nonzero, by
Nakayama's lemma, and $D$ is a finite-dimensional division algebra
over $\mathbb Q$.  The target therefore has rational dimension
$\dim_{\mathbb Q}D$ times the product of the two positive $D$-ranks.
Consequently $D=\mathbb Q$, both ranks are one, and the displayed
surjection is an isomorphism.  Every balanced functional on
$W\otimes_A V$ thus factors through the two one-dimensional tops.
Its bilinear matrix has rank at most one.  The square vanishes by
Lemma~\ref{B:rank-one-square}.
\end{proof}

\begin{proposition}\label{B:local-three-by-three-product}
Let $A$ be a five-dimensional local $k$-algebra with residue field $k$.
There are no right and left modules $W,V$, both of dimension three, for
which
\[
 \dim_k\operatorname{Tor}^A_i(W,V)=
 \begin{cases}2&i=0,\\1&i=2,\\0&i=1\text{ or }i\geq3.\end{cases}
\]
More generally, it suffices to assume the displayed dimensions for
$0\leq i\leq3$ and eventual vanishing of the higher Tor groups.
\end{proposition}
\begin{proof}
Neither module is projective, since a nonzero free $A$-module has
dimension divisible by five.  Lemma~\ref{B:top-heavy-Tor}, applied on
either side, shows that both modules have one-dimensional top and
two-dimensional radical.  In particular $V$ is cyclic.  Its first
syzygy has dimension two.  If this syzygy were not cyclic, its top
would have dimension two and its radical would be zero; dimension
shifting and Lemma~\ref{B:top-heavy-Tor} would make $W$ projective.
Thus the first syzygy is cyclic and its minimal projective cover is
$A$.  Its next syzygy has dimension three and must again be cyclic:
otherwise its top has dimension at least two and its radical has
dimension at most one, giving the same contradiction.  Repeating
this argument shows that every term in the minimal resolution of
$V$ is one copy of $A$, while its syzygies alternate in dimension
between three and two.

Tensor this resolution with $W$ and let $r_i$ be the rank of the
differential from degree $i$ to degree $i-1$.  All entries of the
minimal differentials lie in $J$, so $r_i\leq\dim WJ=2$.  The asserted
dimensions in degrees zero through three force, successively,
\[
 r_1=1,\qquad r_2=2,\qquad r_3=0,\qquad r_4=3.
\]
The last equality contradicts $r_4\leq2$.
\end{proof}

\begin{remark}\label{B:local-three-by-three-scope}
For a finite monoid whose minimal ideal is a $4\times4$ rectangular
band and which has five elements outside that ideal, the two
augmentation modules in the band Tor formula both have dimension
three.  Proposition~\ref{B:local-three-by-three-product} therefore
excludes a local contracted quotient for a classifying space with
the rational homology of $S^2\times S^2$.  The proposition makes no
corresponding assertion for the $\mathbb{CP}^2$ profile: in that case
the same rank equations are $2,1,1,2$, which do not by themselves
give a contradiction.
\end{remark}
